\documentclass{article}
\usepackage{tikz-timing}

\begin{document}

\section{Timing protocols}
The smart cubes communicate using a custom protocol over two data lines.
Each cube has two data lines, D1 and D2.
The data lines are only used for communication between cubes.
Note that the cube's data lines are inverted, meaning D1 of cube A is connected to D2 of cube B and vice versa.

\subsection{Announcement}
The announcement is used to connect a new cube inside a powered grid of cubes.
The new cube will announce its presence to one side, starting at the first side (top).
If another cube is present at that side, it will acknowledge the announcement of the new cube and a handshake will be performed.
During the handshake, both cubes will check both their data lines for connection.

If presence is not acknowledged or the handshake fails, the new cube will retry the announcement on the next side.

The following timing diagram shows the digital signals.
The Bus in the following diagram represents the actual state of the data lines of cube A.
The data lines of cube A and cube B show their respective internal signals.
A state of high impedance implies that the pin for that cube is set as input.

\vspace{10pt}
\begin{tikztimingtable}[scale=3, xscale=1.2]
Bus 1 & HH;[dotted]2H;[green]LL;H;[red]L;
S@{\draw[red,line width=0.8pt]   (\xunit*8,  0pt) -- (\xunit*9,  1pt);
   \draw[green,line width=0.8pt] (\xunit*8, -1pt) -- (\xunit*9,0pt);}
;[green]L \\
Bus 2 & H;[red]L;[red][dotted]2L;[red]L;H;[green]LL;H;[red]L \\
Cube A D1  & [green]HH;[dotted]2HN[rectangle,yshift=5pt,black]{ack presence};LLZZLL \\
Cube A D2  & [green]ZZ;[dotted]2Z;ZZN[rectangle,yshift=5pt,black]{ack handshake};LLZZ \\
Cube B D2  & [red]ZZ;[dotted]2Z;ZZZN[rectangle,yshift=5pt,black]{complete handshake}LLZ \\
Cube B D1  & [red]HN[rectangle,yshift=5pt,black]{announce}L;[dotted]2L;LN[rectangle,yshift=-5pt,black]{init handshake}ZZZZL \\
Note  & SS2D{delay \(\le\) 10 ms} \\
\end{tikztimingtable}

\subsection{Data transfer (initiation)}
When two cubes are successfully connected, they are in the idle state.
Either cube can initiate data transfer to the other cube by asking for a communication.
If the data initiation process is successful, the cube that asked for communication can start sending data.

The first byte of data sent is always the length of the data packet in bytes.
The next bytes are the actual data bytes.

If the data specifies that a response is expected, the receiving cube will send a response packet back to the sender.
Else the communication is finished and the cubes return to idle state.

The following timing diagram shows the digital signals.
The Bus in the following diagram represents the actual state of the data lines of cube A.
The data lines of cube A and cube B show their respective internal signals.
A state of high impedance implies that the pin for that cube is set as input.

\vspace{10pt}
\begin{tikztimingtable}[scale=3, xscale=1]
Bus 1 & ;[green]L
S@{\draw[green,line width=0.8pt]   (\xunit*1,  0pt) -- (\xunit*2,  1pt);
   \draw[red,line width=0.8pt] (\xunit*1, -1pt) -- (\xunit*2,0pt);}
;[red]LL;HH;[red]2D{[black]data};H;[green]2D{[black]data};[green]L \\
Bus 2 & ;[red]L;HH;[green]LL;H;[red]2D{[black]data};H;[green]2D{[black]data};[red]L \\
Cube A D1 & [green]LLN[rectangle,yshift=-5pt,black]{check conn.}ZZZZZZH;[green]2D{[black]data}L \\
Cube A D2 & [green]ZZZN[rectangle,yshift=5pt,black]{ack}LLN[rectangle,yshift=-5pt,black]{wait for data}ZZZH;[green]2D{[black]data};Z\\
Cube B D2 & [red]ZN[rectangle,yshift=5pt,black]{ask}LLLN[rectangle,yshift=-5pt,black]{ack received}HH;[red]2D{[black]data};ZZZZ\\
Cube B D1 & [red]LZZZZZ;[red]2D{[black]data};ZZZL\\
\end{tikztimingtable}

\subsection{Data transfer}
When the data transfer initiation succeeds, data can be transferred.
This process uses one line as a clock and the other line for sending data bit by bit.

The following timing diagram shows the digital signals for a short example data transfer.
The bits are read on falling edges of the clock.
Therefore, the data line must be stable before the falling edge of the clock and can change after the rising edge.

\vspace{10pt}
\begin{tikztimingtable}[scale=2, xscale=1]
Bus 1 (clock) & HCCCCCCCCCCCCCCCC \\
Bus 2 (data) & 0.5H2H2L2L2H2L2H2H2H0.5H \\
Data bits & [d]2D{1}2D{0}2D{0}2D{1}2D{0}2D{1}2D{1}2D{1}0.2D \\
\end{tikztimingtable}

\end{document}
